% sections/methods.tex -- Methods: computation of the series, numerical approach, validation strategy
\section{Methods}
\label{sec:methods}

This section explains how the series coefficients are reduced to integrals (Sec.~\ref{sec:methods-series}), how
the graphs and their integrands are built (Sec.~\ref{sec:methods-graphs}), how the integrals are evaluated
(Sec.~\ref{sec:methods-numerics}), how the series are resummed (Sec.~\ref{sec:methods-resummation}), and how
the results are checked (Sec.~\ref{sec:methods-validation}). A claim that a script of the repository verifies is marked with the identifier of
its check, as in \check{G2}; Appendix~\ref{app:reproducibility} lists the checks.

% ---------------------------------------------------------------------------------------------------------------
\subsection{Computing the series}
\label{sec:methods-series}

\subsubsection{Scheme and Callan--Symanzik relations}
\label{sec:methods-cs}

We consider the Euclidean action
\begin{equation}
S[\phi]=\int d^3x\,\Bigl[\tfrac12(\partial_\mu\phi)^2+\tfrac12 m_0^2\,\phi^2+\tfrac{\lambda_0}{4!}\,(\phi^2)^2\Bigr]
\label{eq:action}
\end{equation}
for an $N$-component field $\phi$, and denote by $\Gamma^{(a,b)}$ the one-particle-irreducible (1PI) vertex functions
with $a$ legs and $b$ insertions of $\tfrac12\phi^2$ at zero momentum, with the $O(N)$ tensors factored
out~\cite{ZinnJustinBook,Ribeca2003thesis}. In the massive zero-momentum scheme~\cite{Parisi1980,BNGM1976,BNM1978}
the renormalized mass $m$, the normalizations of $\phi$ and $\phi^2$, and the dimensionless coupling $g$ are fixed by
\begin{gather}
\Gamma^{(2)}_R(p)=m^2+p^2+O(p^4),\nonumber\\
\Gamma^{(4)}_R(0)=m\,g,\qquad \Gamma^{(2,1)}_R(0;0)=1 .
\label{eq:normcond}
\end{gather}
The renormalized functions obey the Callan--Symanzik equation~\cite{ZinnJustinBook,Ribeca2003thesis}
\begin{equation}
\Bigl[m\,\partial_m+\beta\,\partial_g-\tfrac a2\,\eta-b\,(\eta_{\phi^2}-\eta)\Bigr]\Gamma^{(a,b)}_R
=(2-\eta)\,m^2\,\Gamma^{(a,b+1)}_R
\label{eq:cs}
\end{equation}
for $b\le1$, where the right-hand side carries one more insertion at zero momentum, and $\beta=m\,\partial g/\partial m$,
$\eta=m\,\partial\ln Z/\partial m$ and $\eta_{\phi^2}=m\,\partial\ln Z_{\phi^2}/\partial m$ are taken at fixed bare
parameters ($Z$ and $Z_{\phi^2}$ normalize $\phi$ and $\phi^2$). Applied at zero momenta to the conditions~\eqref{eq:normcond}, the first through its
$p^2$ derivative, Eq.~\eqref{eq:cs} gives, with $u=g/8\pi$ and $\beta_u=m\,\partial u/\partial m$,
\begin{gather}
\eta=\frac{2h}{1+h},\qquad \beta_u=-u+2\eta\,u+(2-\eta)\,\tilde b,\nonumber\\
\eta_{\phi^2}=-(2-\eta)\,c,\qquad \frac1\nu=(2-\eta)(1-c),
\label{eq:csrel}
\end{gather}
where
\begin{gather}
\tilde b=\frac{m\,\Gamma^{(4,1)}_R(0)}{8\pi},\qquad c=m^2\,\Gamma^{(2,2)}_R(0),\nonumber\\
h=-m^2\,\frac{\partial\Gamma^{(2,1)}_R}{\partial p^2}\bigg|_{p=0}.
\label{eq:bhc}
\end{gather}
The exponents follow at the non-trivial zero of $\beta_u$ (Sec.~\ref{sec:methods-resummation}).

The three functions~\eqref{eq:bhc} are renormalized vertex functions with insertions, superficially convergent, and
their expansions in $u$,
\begin{gather}
\tilde b=\sum_{L\ge1}(-1)^{L+1}B_L\,u^{L+1},\nonumber\\
h=\sum_{L\ge2}(-1)^L H_L\,u^L,\qquad c=\sum_{L\ge1}(-1)^{L+1}C_L\,u^L,
\label{eq:series}
\end{gather}
are \emph{linear} in the graphs: $B_L$ is a weighted sum over the $L$-loop graphs of $\Gamma^{(4)}$ with one insertion,
$H_L$ and $C_L$ over the $L$-loop graphs of $\Gamma^{(2)}$ with one and two insertions, of integrals in which every
renormalization part is subtracted at zero momentum (Secs.~\ref{sec:methods-param} and~\ref{sec:methods-rbar}). No $Z$ factor
appears, because the subtractions implement the conditions~\eqref{eq:normcond} graph by graph. Together with lower orders, the coefficient of
$\ubar^{L+1}$ in $\beta$ is fixed by $B_L$ and, through the term $2\eta u$, by $H_L$; that of $\ubar^{L}$ in $\eta$
by $H_L$; and that of $\ubar^{L}$ in $\nu$ by $C_L$ and $H_L$ [Eq.~\eqref{eq:linear} in
Appendix~\ref{app:conventions}]. Only products and quotients of series enter Eq.~\eqref{eq:csrel}; no series has to
be reversed. At lowest order, for instance, $B_1=(N+8)/12$ gives $\beta(\ubar)=-\ubar+\ubar^2+O(\ubar^3)$ for every
$N$. We write the RG functions in the normalization of Ref.~\cite{Phi4tools2024},
\begin{align}
\beta(\ubar)&=-\ubar+\ubar^2+\sum_{k\ge3}\frac{p_{\beta,k}(N)}{(N+8)^{k-1}}\,\ubar^{k},\nonumber\\
\eta(\ubar)&=\sum_{k\ge2}\frac{p_{\eta,k}(N)}{(N+8)^{k}}\,\ubar^{k},\label{eq:series-form}\\
\eta_{\phi^2}(\ubar)&=\sum_{k\ge1}\frac{p_{\eta_{\phi^2},k}(N)}{(N+8)^{k}}\,\ubar^{k},\nonumber\\
\nu(\ubar)&=\frac12+\sum_{k\ge1}\frac{p_{\nu,k}(N)}{(N+8)^{k}}\,\ubar^{k},\nonumber
\end{align}
with polynomials $p_{X,k}(N)$, which follow from the polynomials $B_L(N)$, $H_L(N)$, $C_L(N)$ through
Eq.~\eqref{eq:csrel} applied to the polynomials themselves, in exact rational arithmetic, rather than to their values
at integer $N$, whose interpolation is ill-conditioned.
Equation~\eqref{eq:csrel} is the fixed-dimension counterpart of the representation of the minimal-subtraction RG
functions by finite subtracted integrals~\cite{AdzhemyanKompaniets2011,AdzhemyanKompaniets2014}. We call this computation, through
Eqs.~\eqref{eq:csrel}--\eqref{eq:series}, route A.

The previous complete computations~\cite{NMB1977,BNM1978,Phi4tools2024} and the automated approach of
Ref.~\cite{Ribeca2003thesis} instead evaluated the zero-momentum vertex functions of a scheme in which only two-point
subgraphs are subtracted (in Refs.~\cite{Ribeca2003thesis,Phi4tools2024} only the tadpole, set to zero, and the sunset,
subtracted at $p^2=0$; in Ref.~\cite{NMB1977}, according to Ref.~\cite{Phi4tools2024}, all two-point subgraphs),
and obtained the RG functions through the $Z$ factors, the relation between the coupling of that scheme and the
renormalized one, and its inverse~\cite{Ribeca2003thesis}; we use this $Z$-factor route as an independent check
(route B, Sec.~\ref{sec:methods-validation}). Along that route, in the intermediate scheme of
Refs.~\cite{Ribeca2003thesis,Phi4tools2024}, the contribution of the $L$-loop graphs is a small difference of large numbers: the $L$-loop four-point graphs contribute
to the coefficient of $\beta$ an amount larger than the net $B_L$ by factors of 5.5, 14.4, 29.2, 49.5 and 74 at
$L=2,\dots,6$, and the $L$-loop graphs of $\partial\Gamma^{(2)}/\partial p^2$ contribute to $\eta$ an amount larger than
$H_L$ by about 20, 114 and 106 at $L=4$, 5 and 6; the rest is cancelled by products of lower orders. The relative
error of the coefficient is the relative error of the graph sum multiplied by this factor, so that the number of samples
grows with its square. In Eq.~\eqref{eq:series} the same cancellations take place within each integrand, between a
graph and its counterterms evaluated at the same point, so that they cost no additional samples;
Sec.~\ref{sec:results-performance} quantifies the gain.

\subsubsection{Parametric integrands}
\label{sec:methods-param}

At zero external momenta and unit mass, the integral of a graph $G$ with $\ell$ lines and $L$ loops, with one
$d^3k/\pi^2$ per loop (Nickel's normalization~\cite{NMB1977}), is
\begin{equation}
A_G=\frac{\Gamma(\ell-3L/2)}{\pi^{L/2}}\int_\Delta d\mu\;U_G^{-3/2},\quad U_G=\sum_{T}\prod_{e\notin T}x_e ,
\label{eq:param}
\end{equation}
where $T$ runs over the spanning trees of $G$, $\Delta$ is the simplex $x_e\ge0$, $\sum_e x_e=1$, and
$d\mu=\delta(1-\sum_e x_e)\prod_e dx_e$~\cite{BergereZuber1974,Ribeca2003thesis,BrownKreimer2013}; the one-loop bubble, for instance,
equals one. The Kirchhoff polynomial $U_G$ factorizes over the 2-connected components (blocks) of $G$. An insertion of
$\tfrac12\phi^2$ at zero momentum on line $e$ squares its propagator, that is, it multiplies the Schwinger integrand by
the parameter of $e$: in Eq.~\eqref{eq:param} it adds a factor $x_e$ and raises the argument of the Gamma function by
one. For a two-point graph with external vertices $a$ and $b$, the external momentum enters as
$\exp(-p^2\,W_G/U_G)$, where $W_G$ is the Kirchhoff polynomial of the graph with $a$ and $b$ identified (equivalently,
the polynomial of the spanning 2-forests that separate $a$ from $b$; $W_G=0$ if $a=b$). The $p^2$ derivative at
$p=0$ is therefore $-\Gamma(\ell-3L/2+1)\,\pi^{-L/2}\int_\Delta d\mu\,W_G\,U_G^{-5/2}$, a single one-signed integrand,
which replaces the expansion into $n(n-1)/2+2n$ signed terms per chain of $n$ propagators used in momentum
space~\cite{Ribeca2003thesis}.

Since $\ell=2L$ for $\Gamma^{(4)}$ and $\ell=2L-1$ for $\Gamma^{(2)}$, all integrals entering Eq.~\eqref{eq:series} carry the
same prefactor $\Pi_L=\Gamma(L/2+1)\,\pi^{-L/2}$, and the per-order sums are
\begin{align}
B_L&=\Pi_L\sum_{G\in\mathcal G_4(L)}S_G\sum_{e}\int_\Delta d\mu\;x_e\,
      \Rbar_{G,e}\bigl[U_G^{-3/2}\bigr],\nonumber\\
H_L&=\Pi_L\sum_{G\in\mathcal G_2(L)}S_G\sum_{e}\int_\Delta d\mu\;x_e\,
      \Rbar_{G,e}\bigl[W_G\,U_G^{-5/2}\bigr],\label{eq:BHC}\\
C_L&=\Pi_L\sum_{G\in\mathcal G_2(L)}S_G\sum_{(e,e')}\int_\Delta d\mu\;x_e\,x_{e'}\,
      \Rbar_{G,ee'}\bigl[U_G^{-3/2}\bigr],\nonumber
\end{align}
where $\mathcal G_E(L)$ is the list of $L$-loop graphs with $E$ legs, $S_G$ the weight of $G$ including its $O(N)$
factor (Sec.~\ref{sec:methods-graphs}), $e$ runs over the lines and $(e,e')$ over the ordered pairs of lines,
$e=e'$ included, and $\Rbar_{G,e}$ is the subtraction operator of Sec.~\ref{sec:methods-rbar}, which depends on where the
insertions are. The propagators have the renormalized mass, $m=1$. (Refs.~\cite{NMB1977,Ribeca2003thesis,Phi4tools2024}
use propagators of unit mass in the scheme of their integrals instead; individual integrals differ, the RG functions
do not.) The sums are known exactly at
one and two loops [Eq.~\eqref{eq:exactsums}]. On the simplex the integrands can be
multiplied by $(\sum_e x_e)^{-(L/2+1)}$, which equals one there and makes them homogeneous of degree $-\ell$: the
integrals are then projective, and can be evaluated in any coordinates of the projective simplex
(Sec.~\ref{sec:methods-tropical}). The integration dimension is $\ell-1$, that is 15 for the eight-loop four-point and 16
for the nine-loop two-point graphs.

\subsubsection{Zero-momentum subtractions}
\label{sec:methods-rbar}

The renormalization parts of a graph are its 1PI subgraphs $\gamma$ of three types, classified by their legs (the
lines and external legs of $G$ attached to $\gamma$) and by the insertions on their lines: four-point parts without
insertions, subtracted by $t^0_\gamma$, the term of order zero of the Taylor expansion in the external momenta of
$\gamma$; two-point parts without insertions, subtracted by $t^2_\gamma$, the terms up to order two (mass and field
normalization); and two-point parts with one insertion, subtracted by $t^0_\gamma$. These subtractions enforce the
conditions~\eqref{eq:normcond} graph by graph~\cite{BogoliubovParasiuk1957,Hepp1966,Zimmermann1969,BergereLam1976,Zavialov1990}. In $d=3$ only the tadpole and
the sunset subgraphs need renormalization~\cite{Ribeca2003thesis,Phi4tools2024}; all the other subtractions are finite, but the scheme requires them,
and through them the cancellations described above take place inside each integrand. The graphs of Eq.~\eqref{eq:BHC}
are themselves not renormalization parts, except those of $H_L$, whose own subtraction $t^0$ does not depend on $p^2$
and drops out of the derivative, so that the subtraction of the subgraphs alone, $\Rbar$, is the complete
renormalization:
\begin{equation}
\Rbar_G=\sum_{F}\prod_{\gamma\in F}\bigl(-t_\gamma\bigr),
\label{eq:forest}
\end{equation}
where $F$ runs over the forests, sets of renormalization parts that are pairwise nested or vertex-disjoint (the
compatibility condition of the counterterms generated by a Lagrangian), the empty forest included, and inner operators
act first.

In parametric space the subtractions have closed forms. The term $t^0_\gamma$ replaces $U_\theta$, in the parent $\theta$ of
$\gamma$, by $U_\gamma\,U_{\theta/\gamma}$, where $\theta/\gamma$ is $\theta$ with $\gamma$ contracted to a point: this is the
leading term of $U_\theta$ when the parameters of $\gamma$ go to zero. For a two-point part $\gamma$ joined to the rest at
the vertices $a$ and $b$, the classical factorization of the Kirchhoff polynomial across a two-vertex
cut~\cite{Nakanishi1971,BognerWeinzierl2010},
\begin{align}
U_\theta&=U_\gamma\,U_{\theta/\gamma}+W_\gamma\,U_{\theta\setminus\gamma}\nonumber\\
&=U_\gamma\,\bigl(U_{\theta/\gamma}+\varepsilon_\gamma\,U_{\theta\setminus\gamma}\bigr),\qquad
\varepsilon_\gamma=\frac{W_\gamma}{U_\gamma},
\label{eq:effline}
\end{align}
where $\theta\setminus\gamma$ is $\theta$ without the lines of $\gamma$ and $W_\gamma$ is the Kirchhoff polynomial of $\gamma$
with $a$ and $b$ identified (a spanning tree of $\theta$ contains either a spanning tree of $\gamma$, and then a
spanning tree of $\theta/\gamma$, or a spanning two-forest of $\gamma$ that separates $a$ from $b$, and then its other
lines connect $a$ and $b$ through $\theta\setminus\gamma$), shows that $\gamma$ acts on its parent as a massless line
with Schwinger parameter
$\varepsilon_\gamma$; $t^2_\gamma$ keeps the terms of order zero and one in $\varepsilon_\gamma$, and reduces to
$t^0_\gamma$ when $a=b$. The term of a forest $F$ is thus
\begin{equation}
T_F=(-1)^{|F|}\,x_{\mathrm{ins}}\prod_{\theta\in F\cup\{G\}}U_{\theta/F}^{-3/2},
\label{eq:forestterm}
\end{equation}
where $x_{\mathrm{ins}}$ is the monomial of the insertions and $\theta/F$ is the minor of $\theta$ in which every maximal
element of $F$ below $\theta$ is contracted ($t^0$) or replaced by its effective line ($t^2$), with the result expanded to
first order in every $\varepsilon$; for $H_L$ the factor of the root is $W_{G/F}\,U_{G/F}^{-5/2}$. Both terms of
Eq.~\eqref{eq:effline} have the degree of $U_\theta$, so all forest terms are homogeneous of the same degree: they share
the simplex and the prefactor of Eq.~\eqref{eq:BHC}, and their sum, with all positions of the insertions, is integrated
as one function~\cite{BrownKreimer2013}.

% sections/graphs.tex -- Methods B: graphs and integrands (generation, construction, consistency checks)
% ---------------------------------------------------------------------------------------------------------------
\subsection{Graphs and integrands}
\label{sec:methods-graphs}

The pipeline from the Lagrangian to the sampled integrands has three stages, each with its own checks: for each order, the generation
of the graphs, the construction of the subtracted integrand of every graph, and the consistency checks that
tie the two to exact results. All three are automatic and were run here to ten loops.

\subsubsection{Generating the graphs}
\label{sec:methods-generation}

The lists of 1PI graphs of $\Gamma^{(2)}$ and $\Gamma^{(4)}$ are produced by a generator that derives from the one
written for the seven-loop computation of Ref.~\cite{Ribeca2003thesis}, with three improvements: it grows the graphs
by vertex insertion instead of enumerating labelled graphs and screening them by a canonical form, it canonicalizes
them by colour refinement with backtracking, and it takes any list of vertex degrees, the external legs being implicit
in the valences. The class it generates is every connected graph with the given valences, self-loops and bridges
allowed, and the construction is complete: removing a leaf of a spanning tree and re-pairing the half-lines it leaves
is the inverse of an insertion, so that every graph with $V$ vertices descends from one with $V-1$. The class
is grown order by order, each candidate canonicalized and kept in a hash set, and the 1PI graphs are written out at
the end of every order, together with their Nickel indices and weights. The cost is a few hundred insertions per
graph grown, so that it follows the size of the class, growing by a factor of ten to twelve per loop on one core.
Table~\ref{tab:generation} gives the sizes and the times, measured on one core of a cluster node; the
counts through eight loops are those of Ref.~\cite{Ribeca2003thesis} (Table~3.1), the others are new~\check{G1}. Every
complete list satisfies its zero-dimensional sum rule (Sec.~\ref{sec:methods-checks}) through twelve loops for
$\Gamma^{(2)}$ and eleven for $\Gamma^{(4)}$.

Every graph is converted to its Nickel index~\cite{NMB1977} in canonical form---the
lexicographic minimum over vertex labellings---which identifies it across lists and keys its random numbers
(Sec.~\ref{sec:methods-rng}); one-vertex-reducible graphs and graphs with both legs on one vertex are flagged, since
their integrands vanish (Sec.~\ref{sec:methods-construction}), and are not built for $B_L$ and $H_L$.

\begin{table}
\caption{\label{tab:generation}%
Generation of the graph lists: numbers of 1PI graphs without self-loops (one-vertex-reducible ones included), and the
CPU time of the generator on one core of a cluster node (an Intel Xeon Platinum 9242), one list at a time on an
otherwise idle node, the Nickel indices, weights and sum rules included. The eleven- and twelve-loop lists were made
on 96 threads, in 23~s and 275~s of wall time for $\Gamma^{(2)}$ and 1005~s for $\Gamma^{(4)}$ at eleven loops, the
last with 38~GB of memory; their times in the table are core times summed over the threads. Every regenerated list
equals the stored one. The counts of graphs whose integrands survive the pruning (Sec.~\ref{sec:methods-construction})
are in Table~\ref{tab:graphcounts}.}
\begin{ruledtabular}
\setlength{\tabcolsep}{3pt}\footnotesize
\begin{tabular}{lrrrr}
$L$ & $\Gamma^{(4)}$ & time & $\Gamma^{(2)}$ & time \\
\hline
7  & 3986       & 2.1 s & 317    & 0.1 s \\
8  & 26\,540    & 22 s & 1622   & 1.1 s \\
9  & 197\,022   & 4.3 min & 9184   & 9.8 s \\
10 & 1\,593\,254 & 53 min & 59\,614 & 98 s \\
11 & 13\,975\,511 & 14.1 h & 428\,164 & 22 min \\
12 & --- & --- & 3\,393\,754 & 4.2 h \\
\end{tabular}
\end{ruledtabular}
\end{table}

\paragraph{Weights.}
Each graph carries the weight $S_G=T_G(N)\,E!/|\mathrm{Aut}\,G|$, where $E$ is its number of legs, $\mathrm{Aut}\,G$ its
automorphism group at the level of half-edges with unlabelled legs, and $T_G(N)$ its $O(N)$ factor, obtained by
contracting the vertex tensors $(\delta_{ij}\delta_{kl}+\delta_{ik}\delta_{jl}+\delta_{il}\delta_{jk})/3$ over the $3^V$
pairings at the $V$ vertices and projecting on the external tensor, so that $T_G(1)=1$. The factor $E!/|\mathrm{Aut}\,G|$
includes the sum over the labellings of the legs, which is legitimate at zero external
momenta~\cite{Ribeca2003thesis}. The integrals do not depend on $N$: since every graph carries its polynomial
$T_G(N)$, one run gives the coefficients for all $N$.

\subsubsection{Constructing the integrands}
\label{sec:methods-construction}

Every polynomial in Eq.~\eqref{eq:forestterm} is the Kirchhoff polynomial of a minor of $G$, possibly with two vertices
identified, and factorizes into the Kirchhoff polynomials of its blocks, which we call atoms. Expanding all forest terms
into products of half-integer powers of atoms, with binomial coefficients that are dyadic rationals and hence exact in
binary floating point, and adding the coefficients of identical products, removes every combination of forest terms
that cancels identically; we call this step the pruning. Each graph becomes
\begin{equation}
f_G(x)=\sum_k c_k\,x_{\mathrm{ins},k}\prod_a P_a(x)^{p_{ka}/2},
\label{eq:integrand}
\end{equation}
with atoms $P_a$, integer $p_{ka}$ and exact $c_k$; terms with the same atoms, which differ only by the insertions, are
grouped. The result of the expansion is a small program: the list of atoms (each a minor of $G$, evaluated at run time
by a spanning-tree pass and a factorization of its loop matrix), the list of terms (coefficient, insertion monomial,
atom exponents), and the tables of the sampling density of Sec.~\ref{sec:methods-tropical}. One generic kernel
evaluates any such program (Sec.~\ref{sec:methods-kernel}).

The pruning removes up to two thirds of the raw terms. At six loops a four-point graph has on average 250 forests,
which leave 116 terms in 22 groups with 19 distinct atoms (at most 55); at seven loops, in a sample of 300 four-point
graphs, 583 forests, 219 terms, 39 groups and 26 atoms (at most 6128, 766, 126 and 85). The largest programs are
those of two-point graphs with many nested two-point parts: at nine loops the chain of eight bubbles of $C_9$ has
$3.4\times10^5$ forests and 2328 atoms, against at most 269 atoms for $H_8$ and 203 for $B_8$. The expansion is done
once per graph on the host, in parallel over the graphs; its cost grows faster than the number of graphs
(Table~\ref{tab:graphcounts}: about \buildpc{b8} core-seconds per graph for $B_8$, \buildpc{h9} for $H_9$,
\buildpc{c9} for $C_9$, and from \buildpc{b10} for $B_{10}$ to \buildpc{c10} for $C_{10}$ at ten loops), so that at ten
loops it dominates a job. The built integrands of a list are therefore
written once to a cache file, by a CPU-only job, and loaded by every subsequent run; a run with and without the cache
gives bitwise the same results~\check{I4}.

Graphs that vanish identically come out of the pruning with no terms at all, without any structural rule: graphs with a
cut vertex (one-vertex-reducible graphs), which factorize at zero momentum~\cite{Ribeca2003thesis} into pieces of which
one is a renormalization part at exactly zero momentum, cancelled by its own subtraction; graphs with self-loops,
for $L\ge2$; and, for $H_L$, generalized tadpoles, graphs with both external legs on one vertex, for which $W_G=0$. At
three to seven loops these are 21--30\% of the four-point and 15--17\% of the two-point lists (more for $H_L$); at
eight loops 18\% and 14\%, and at nine loops 12\% of the two-point list. In production runs the one-vertex-reducible
graphs of $B_L$ and $H_L$ are dropped before the expansion; their vanishing, and that of the graphs with self-loops, was
verified explicitly up to five loops~\check{I2}. Those of $C_L$ are expanded in every run, and the pruning leaves no
term in any of them, at every order (independently verified up to seven loops~\check{I2}).

\paragraph{Twins.}
At zero momenta and equal masses, Eq.~\eqref{eq:param} depends on $G$ only through its cycle matroid, so that
2-isomorphic graphs in Whitney's sense~\cite{Whitney1933}, which we call twins, have equal integrals in any
representation~\cite{Ribeca2003thesis}: the 3986 seven-loop four-point graphs have 3018 distinct unsubtracted integrals,
and the 3150 non-vanishing ones fall into 2391 classes. We checked that the subtracted integrands of twins, summed over
the positions of the insertion, agree pointwise under the matroid map of the lines (all 3 pairs at four loops and all 25
non-vanishing pairs at five loops for $B_L$)~\check{I3}. Twins have different $O(N)$ factors, so that a merged class carries the
sum of their weights. Twins were not merged in the runs of Sec.~\ref{sec:results}: every class is one graph.

\subsubsection{Consistency checks}
\label{sec:methods-checks}

The graphs, their weights and their integrands are checked, before any sampling, against results that are exact.
\begin{enumerate}
\item \emph{Zero-dimensional sum rules.} In zero dimensions every integral equals one, so that $\sum_{G}S_G$ over a
complete list is the coefficient of $g^L$ in the corresponding vertex function of the zero-dimensional $O(N)$ model,
known exactly from its moments (Appendix~\ref{app:sumrules}). Summed over all graphs, one-vertex-reducible ones
included, this tests the completeness of a list and its weights without any integral; for lists without self-loops the
bare mass is shifted so that every self-loop cancels. Every list passes it exactly, with the full $O(N)$
polynomial, through ten loops (Table~\ref{tab:sumrules})~\check{G2}.
\item \emph{Independent lists.} Converted to Nickel indices, our lists equal those of Ref.~\cite{Phi4tools2024}
(generated with \textsc{feyngen}~\cite{Borinsky2014feyngen}) as sets, for $\Gamma^{(4)}$ at $L\le7$ and $\Gamma^{(2)}$ at
$L\le8$~\check{G3}, and those of a brute-force generator in our code where that is practical ($\Gamma^{(4)}$ to five loops,
$\Gamma^{(2)}$ to six). Two independent canonicalizers (one in the generator's string format, one on Nickel indices)
agree on every graph, and all graphs of every list are distinct~\check{G1}.
\item \emph{Weights.} The weights and $O(N)$ polynomials equal those of the independent implementation of
Refs.~\cite{Ribeca2003thesis,GuidaRibeca2006} on all 260 graphs with $L\le5$ (its 251 and the nine graphs with self-loops its lists lacked) and on all 47
graphs of its four-loop tables~\check{G4}.
\item \emph{Integrands.} The expansion~\eqref{eq:integrand} agrees pointwise with an independent evaluation of the
forest formula in 100-digit arithmetic (Sec.~\ref{sec:methods-validation})~\check{I1}; vanishing graphs
vanish~\check{I2}, twins agree pointwise~\check{I3}, and the exact two-loop sums $B_2=1/2$, $H_2=1/81$ and $C_2=1/12$
are reproduced to a few $10^{-7}$ by the sampled integrals (Sec.~\ref{sec:results-validation})~\check{N3}.
\item \emph{Large order.} Each new coefficient is compared with its leading large-order estimate
(Sec.~\ref{sec:methods-validation}), a test of the complete chain that catches an error of a few per cent (Sec.~\ref{sec:methods-validation})~\check{S3}.
\end{enumerate}


% ---------------------------------------------------------------------------------------------------------------
\subsection{Numerical approach}
\label{sec:methods-numerics}

\subsubsection{Choosing the method}
\label{sec:methods-choice}

The earlier evaluations integrated each graph separately in momentum space, with one-loop subgraphs replaced by
analytic effective vertices, on deterministic meshes~\cite{NMB1977}, by adaptive
cubature~\cite{Ribeca2003thesis,GuidaRibeca2006}, or with VEGAS~\cite{Lepage1978,Phi4tools2024}. Because of the
cancellations of the $Z$-factor route, and of the precision lost in finding the zero of $\beta$ and in the resummation,
eight to ten digits per graph were considered necessary~\cite{Ribeca2003thesis}; high-precision
methods based on integration-by-parts reduction~\cite{ChetyrkinTkachov1981,Laporta2000} are available for massive vacuum integrals only to five
loops~\cite{LutheSchroeder2016}. The cost of that strategy is set by its hardest graphs: every graph must reach many
digits, and both the number of graphs and their dimension grow with the order, so that each new order costs more than
the last by a large and growing factor.

We made the opposite choice deliberately: a method whose precision depends as little as possible on the order. Every
integrand, whatever its size, is sampled by the same algorithm. Like any Monte Carlo estimate, an order sum then
reaches a relative error $\epsilon$ in a time proportional to $1/\epsilon^2$; what the design controls is the constant of
proportionality, and it aims at one that grows as slowly as possible with the order. Three choices serve that aim. The
samples are shared among the graphs of an order in proportion to their spread (Sec.~\ref{sec:methods-estimators}),
so that the constant is set by an average over the graphs of the relative spread and the cost of one sample, and by a
factor that measures the cancellations between graphs of opposite sign: the number of graphs does not enter it by itself. Each
graph is sampled with a density that follows its singularities (Sec.~\ref{sec:methods-tropical}), so that the weights
stay bounded at every order. And each term of the forest formula~\eqref{eq:forestterm} costs only a polynomial in $L$
to evaluate, since its polynomials are determinants of $L\times L$ matrices (Kirchhoff's matrix-tree theorem). At low
orders such a method cannot match the per-graph treatment, where effective vertices reduce each integral to a few
dimensions and deterministic integration reaches eight to ten digits. At high orders, where the dimension grows with
the number of lines and deterministic rules slow down accordingly, it should become far cheaper. How slowly its
constant actually grows could only be measured: empirically, the time to a fixed relative precision grows by a factor of about
\costgrowthrange\ per loop (Sec.~\ref{sec:results-performance}). The
construction of the integrands, done once per graph, still grows with their number, that is factorially. It is a
one-off cost on ordinary CPUs, but at high enough orders it becomes the dominant one: at the tenth order it already
exceeds the sampling (Sec.~\ref{sec:results-performance}). A precision that hardly depends on the order also keeps
the statistical errors of the highest coefficients, which carry the information on the large-order behaviour, from
limiting the resummation; at the present orders they are already negligible there (Sec.~\ref{sec:resummation-errors}).

The lower-order terms that enter a new coefficient together with its new order sums are known far more precisely
and do not cancel them, so that the relative error of a coefficient is that of its order sums. What matters is therefore the relative error of each order sum, which no
cancellation amplifies. Monte Carlo integration then becomes adequate: it requires no choice of variables per
graph, it gives reliable error estimates when the weights are bounded, and it is naturally parallel. The parametric
integrands, however, are singular on every face of $\Delta$ on which the parameters of a subgraph vanish, and the faces
overlap at the corners; under uniform sampling the variance is already infinite on the face of a one-loop bubble. Sector
decomposition~\cite{BinothHeinrich2000,pySecDec2018} removes these singularities, but in our tests it required $10^5$--$3\times10^5$ sectors for a six-loop
graph and $(1$--$5)\times10^6$ for a seven-loop four-point graph, and lattice rules on the sector
integrands~\cite{pySecDecQMC2019} remained pre-asymptotic up to $10^9$ points at 9--11 dimensions. We use instead
tropical sampling~\cite{Borinsky2023tropical,feyntrop2024}, a recent technique, increasingly used for Feynman
integrals~\cite{Balduf2023,feyntrop2024,Borinsky2025global}, which samples the Hepp sectors~\cite{Hepp1966} rather
than enumerating them. A sampler of the same
kind, with an exponent for every subset of lines and the sector drawn by a recursion over the subsets, was used
for the forest-subtracted parametric integrands of the electron anomalous magnetic moment, on CPUs and then on
GPUs~\cite{Volkov2017,Volkov2018,Volkov2024}; our exponents follow the forest terms of the zero-momentum subtraction
(Sec.~\ref{sec:methods-tropical}). Each graph is sampled with its own density. Sampling a whole order with one
density~\cite{Borinsky2025global} brings no gain for the unsubtracted integrands, which are all positive: in our tests at
four and five loops a one-sample global sampler came within 3--4\% of the optimally stratified per-graph sampling
described below. Its advantage lies elsewhere: it draws graphs and points together, without enumerating the graphs, in
polynomial time and memory, and it has been applied to finite correlators of a massive super-renormalizable theory in
three dimensions~\cite{Borinsky2025global}.

\subsubsection{Forest-aware tropical sampling}
\label{sec:methods-tropical}

In the Hepp sector $x_{\sigma(1)}<\dots<x_{\sigma(\ell)}$, with $S_j=\{\sigma(1),\dots,\sigma(j)\}$ and
$t_j=x_{\sigma(j)}/x_{\sigma(j+1)}$, the density with respect to the projective measure $\prod_j dt_j/t_j$ is
\begin{gather}
p(x)=\frac{1}{J(E)}\prod_{j=1}^{\ell-1}t_j^{\,\omega(S_j)},\nonumber\\
J(S)=\sum_{e\in S}\frac{J(S\setminus e)}{\omega(S\setminus e)},\qquad J(\{e\})=1 ,
\label{eq:tropical}
\end{gather}
where $J(E)$, a generalized Hepp bound~\cite{Panzer2022Hepp}, sums $\prod_j\omega(S_j)^{-1}$ over all sectors, and the
exponents $\omega(S)$ follow the integrand. When the parameters of a set $S$ of lines go to zero together,
$x_e\to\lambda x_e$ for $e\in S$, an atom behaves as $\lambda^{h_a(S)}$, where $h_a(S)$ is the loop number of the
lines of $S$ that belong to it, and the term $k$ of Eq.~\eqref{eq:integrand} as $\lambda^{\deg_S k}$, with
$\deg_S k=\sum_a p_{ka}\,h_a(S)/2+n_{k}(S)$ and $n_k(S)$ the number of its insertions on lines of $S$. The term
converges on this face if $|S|+\deg_S k>0$. We take
\begin{equation}
\omega(S)=s\,\Bigl[\,|S|+\min_k\deg_S k+\delta(S)\Bigr],
\label{eq:omega}
\end{equation}
with the minimum over all terms that survive the pruning, $\delta(S)=1$ where $|S|+\min_k\deg_S k\le0$ and zero
elsewhere, and $s\in(0,1]$ a softening factor: $s=1$ gives the tropical density, and $s<1$ lowers every exponent by the
same fraction, which tames the tails of the weights at the price of more samples in the bulk; we use $s=0.7$, chosen as
described in point (iii) below. Both
$\omega$ and $J$ are tabulated on the host for all $2^\ell$ subsets of lines, so that the cost grows as $2^\ell$ rather
than as the number $\ell!$ of sectors; the two tables take 1\,MB per graph at $\ell=16$. A point is drawn by removing, from
$S=E$ downwards, the largest line $e$ of $S$ with probability $J(S\setminus e)/[\omega(S\setminus e)\,J(S)]$, and setting
the ratio of the largest remaining parameter to $x_e$ to $t=v^{1/\omega(S\setminus e)}$ with $v$ uniform~\cite{Borinsky2023tropical};
a sample takes $2(\ell-1)$ uniform numbers. Its weight $w=f_G(x)\prod_e x_e/p(x)$, with $f_G$ in projective form, has expectation
$\int_\Delta d\mu\,f_G$, and the estimator is unbiased for any positive $\omega$.

Three features of Eq.~\eqref{eq:omega} adapt the tropical density to the subtracted integrands. (i) The minimum runs over
all forest terms: a counterterm brings in the faces of its cograph, and a density built from the unsubtracted graph alone
produced weights up to $8\times10^4$ times the mean. (ii) On faces where every term diverges, $|S|+\min_k\deg_S k\le0$, the
sum converges only through a $t^2$ subtraction, which raises the true exponent by two units (for the sunset, $|S|+\deg_S$ goes from 0
to 2, the measure from $\lambda^{-1}d\lambda$ to $\lambda^{+1}d\lambda$) and which the term-by-term minimum does not see; the shift $\delta=1$ keeps the density
integrable and below the true behaviour, as we checked in extended precision on every shifted face of 280 integrands at
four and five loops (no violation); a shift by two units changes the per-sample relative standard deviation of the order sum [$\rho$ of Eq.~\eqref{eq:rho}
below] by less than 0.5\%. On no face is the exponent negative before the shift, for any integrand up to ten loops: the builder refuses such
integrands, on which the double-precision evaluation of the expanded $t^2$ terms would have round-off of infinite variance,
and it built every class of every order~\check{G5}. (iii) The tropical density falls off faster than the integrand before the
asymptotic regime: the positive powers of atoms in the $t^2$ terms give large but bounded prefactors that it does not see,
and it ignores the factor $(\sum_e x_e)^{-(L/2+1)}$, which plausibly explains why softening helps the unsubtracted
integrands as well. With $s=1$ the weight distributions had tails with a Hill index~\cite{Hill1975}
as low as 1.9 and five-loop error estimates were unreliable. Softening all exponents by $s<1$ makes the weights vanish on
every face. The value $s=0.7$ is close to optimal for every kind of integrand at four and five loops, where it halves the
per-sample standard deviation of both the subtracted and the unsubtracted integrands (any positive set of exponents
gives a valid density, which can be fitted to the integrand~\cite{Salvatori2025}; the optimum found there for the
corresponding uniform factor is 0.72); the smallest tail index is then 3.87,
below 4, so that the fourth moment, and with it the error of the error below, is only indicative. We have not repeated
the scan of $s$ and of the shift beyond five loops. At six and seven loops replica ensembles test the resulting error
estimates directly (Sec.~\ref{sec:results-validation})~\check{N7}, and at eight loops the independent runs of $B_8$ and
$H_8$ do (Table~\ref{tab:newsums}); beyond that they rest on the within-run statistics alone.

\subsubsection{Estimators, error estimates and allocation}
\label{sec:methods-estimators}

Every class is sampled independently; a class is one graph, since twins (Sec.~\ref{sec:methods-construction}) were not
merged. Its contribution to the order sum
$X_L=\sum_G a_G\,I_G$, with $a_G=\Pi_L S_G$ and $I_G$ the integral in Eq.~\eqref{eq:BHC} summed over the positions of the insertions, is estimated as
$a_G\bar w_G$ with variance $a_G^2\sigma_G^2/n_G$, where $\bar w_G$ and $\sigma_G$ are the mean and standard deviation of its
$n_G$ weights. For each class we also record the kurtosis $\mathcal K_G$ of the weights, from which the relative standard error
of $\sigma_G$, $[(\mathcal K_G-1)/4n_G]^{1/2}$, is propagated to that of the error of $X_L$ (the error of the error); the largest
weight relative to the mean; and the largest round-off scale of a single sample (Sec.~\ref{sec:methods-rng}). An order is
characterized by
\begin{equation}
\rho=\frac{\sum_G|a_G|\,\sigma_G}{|X_L|},\qquad \rho_\tau=\frac{\sum_G|a_G|\,\sigma_G\sqrt{\tau_G}}{|X_L|},
\label{eq:rho}
\end{equation}
the per-sample relative standard deviation under optimal allocation, such that $n$ samples give a relative error
$\rho/\sqrt n$, and its counterpart per unit of computing time, with $\tau_G$ the measured cost per sample.

After a pilot of $2^{14}$ to $2^{16}$ samples per class in the runs of Sec.~\ref{sec:results}, the samples are allocated by Neyman's
rule~\cite{Neyman1934}, $n_G\propto|a_G|\sigma_G$, or, in a time-boxed
run, $n_G\propto|a_G|\sigma_G/\sqrt{\tau_G}$ with $\tau_G$ in core or device seconds, which minimizes the error at fixed
computing time. The main phase runs in rounds (eight by default), each of which brings every class to a fixed fraction of
its allocation; the complete state is written after every round, and a run that reaches its time limit stops at the end of
a unit of work. Since the samples of a class always form a prefix of its random-number stream, an interrupted run is
identical to a shorter one. At $s=0.7$ the pulls of single runs were Gaussian over 340 replica runs on CPUs at two to four loops,
with no single run beyond three standard deviations, and over 400 replicas at three and four loops with the final code
the mean pulls were within $\pm0.1$ of zero and their spread between 0.91 and 1.06; at $s=1$ fatter tails had been
seen. Since the sample variance of a heavy-tailed class converges slowly, we nevertheless judge agreement over
independent seeds rather than from the error estimate of a single run (Sec.~\ref{sec:methods-validation}).

\subsubsection{Pooling the runs: the ledger}
\label{sec:methods-ledger}
Each class draws its random numbers from its own stream of a counter-based generator (Sec.~\ref{sec:methods-rng}),
whose key is derived from the seed and from the class (kind, loop order, Nickel index). Two runs with different seeds therefore sample every
class independently, whatever their hardware, allocation or duration and whenever they are contributed, and their per-class statistics pool exactly:
with $n_i$ samples, weight mean $m_i$ and weight variance $s_i^2$ from run $i$, the pooled mean is
$\sum_i n_i m_i/\sum_i n_i$ and the pooled sum of squares is $\sum_i [(n_i-1)s_i^2+n_i m_i^2]$. The ledger keeps,
for every class, the registers of all its samples pooled: their number, mean and central moments up to the fourth,
updated run by run with the pairwise formulas of Welford, Chan, Golub and LeVeque, and
P\'ebay~\cite{Welford1962,ChanGolubLeVeque1979,Pebay2008}, together with the extremes of the weights, the round-off
indicators and the time per sample, so that the ledger's size does not grow with the number of runs. For every run it records
the seed, the code version and build signature, the device self-tests, the classes covered and the run's own
statistics; every run file is kept as received in an archive outside the ledger, from which the ledger can be rebuilt
bit for bit~\check{N5}. The per-order sums, the coefficients and their uncertainties, and the tables of this paper are
regenerated from the ledger~\check{S2}. A run is admitted to the ledger only if its seed is new for that
order, its build is a known one whose devices passed the start-up self-test against the host, it has no non-finite
samples, and its classes are statistically consistent with the pooled ones (a $\chi^2$ over the classes it shares with
the ledger, which detects a wrong build or a failing device)~\check{N5}. The ledger and the tools are public with the code; a
refinement is one command that chooses the order with the largest relative uncertainty, allocates the time across its
classes by their measured cost and spread, runs with a fresh seed and submits its contribution. The precision of every
coefficient can thus be improved indefinitely, by anyone, without recomputing what exists.

\subsubsection{Random numbers and arithmetic}
\label{sec:methods-rng}

The random numbers come from the counter-based generator Philox4x32-10~\cite{Salmon2011Random123}, a pure function of a
key and a counter; our implementation reproduces its published known-answer vectors~\check{N1}. Each class has its own key, a hash of
the seed and of a string made of the kind of integrand, the loop order and the Nickel index, and the counter is the sample
index. The random numbers of a class therefore depend only on the seed, the kind, $L$ and the graph, not on its position in
a list, on how an order is split into runs, or on the hardware. Runs at different loop orders are independent even with the
same seed (the largest measured correlation between index-matched classes, $5.4\times10^{-3}$, is consistent with zero), and replicas differ by the seed.

The evaluation is entirely in double precision (FP64), in logarithmic form. For each atom, the leading monomial at the
sampled point is the product of the parameters of the chords of a minimum spanning tree, found by Kruskal's algorithm with
the lines in increasing order of $x_e$; the loop matrix of the corresponding fundamental cycles, rescaled by the chord
parameters, has entries that are ratios of parameters not larger than one and a determinant between one and the number of
spanning trees, computed by an $LDL^{\mathsf T}$ factorization. Hence $\log P_a$ is obtained without overflow or underflow at
any point the sampler can produce, however close to a face. The terms of Eq.~\eqref{eq:integrand} are added with
Neumaier's compensated summation~\cite{Neumaier1974}, and the sums and moments of the weights are accumulated with
compensation as well. Pruning removes the terms that cancel identically; what remains are the cancellations of the
subtractions close to the faces they protect, where double precision loses all relative accuracy pointwise (within
about $3\times10^{-7}$ of a sunset face). The absolute round-off of a sample remains of the order of the machine
precision times $\sum_k|\text{terms}|$, which is negligible against $\sigma_G$: double and 80-bit extended precision at
the same points differ by at most $1.7\times10^{-8}$ standard deviations (rms) at four and five loops, and every run
reports the worst single-sample ratio of this scale to $\sigma_G$, with a warning above $10^{-3}$. Single precision is not
used, as it is not adequate for subtracted integrands: already for the sunset-subtracted five-loop integrals of the
intermediate scheme its error reaches $10^{-2}$ standard deviations per sample.

\subsubsection{A single kernel for CPUs and GPUs}
\label{sec:methods-kernel}

The integrator is written in C++. The sampler and the evaluation of Eq.~\eqref{eq:integrand} are written once, as
header code compiled both by the host C++ compiler and by the CUDA compiler; the expansion of Sec.~\ref{sec:methods-rbar} runs on the host and produces, for
each class, a flat image of its atoms, terms and tables. The unit of computation is a chunk of 4096 consecutive samples
of one class, evaluated by a block of 128 lanes of 32 samples each. Each lane accumulates in registers with compensated
summation, the block reduces its lanes along a fixed binary tree, chunks are merged in index order within units of 16
chunks, and units in order. Compensated merges are not associative bit for bit, so this order is part of the definition
of the result: the accumulators of a class depend only on its key and on its number of samples. On CPUs they are bitwise
identical for any number of threads; an emulator that runs the device program on the host, with the data layout of the
device, reproduces them bitwise for any number of devices and blocks, any launch size and any order of execution. The
kernel body is written against an executor that either loops over the lanes (CPU) or maps them onto threads (GPU), in
phases that end with a barrier by construction; running the lanes in reverse order gives the same bits, which excludes a
missing barrier. A GPU can therefore differ from the CPU only through the device mathematical library and the contraction
of multiplications and additions into fused operations (the host code is compiled without contraction).

This is tested on every machine by golden runs: for 13 configurations (subtracted integrands of all three kinds,
and intermediate-scheme four-point integrands, at three to six loops), reference files computed on a CPU with fixed seeds
and 65\,536 samples per class must be reproduced on the GPUs with identical sample counts and per-class means within 100
units of the round-off of the per-sample terms and within $10^{-4}$ standard deviations, on one device, on each device
separately and with the work forcibly split among all of them~\check{N2}. On the CPU, perturbing every exponential and logarithm by up
to 2 units in the last place moves the per-class means by less than one such unit. Before production each GPU also passes
a self-test: device exponential and logarithm within 2 units in the last place of the host, exact square root and fused
multiply-add, the Philox known answers, a signature of all compile-time limits and data layouts, 4096 samples of each of up to
three classes compared one by one with the host (same Hepp sector, same weight to $10^{-11}$ of the sum of the absolute values of the
terms), and one chunk through the production kernel.

On a GPU the kernel runs as persistent blocks that take chunks from an atomic counter, one class per block at a time, so
that the threads of a warp follow the same integrand and chunks of classes with different costs balance automatically.
Every GPU is driven by its own host thread, which takes units of 65\,536 samples from a common queue in class order and
keeps the images of the classes it uses in device memory, with least-recently-used eviction. Launches hold about half a
second of work but at least two full waves of the grid; their measured durations are apportioned to the classes by their
cycle counts and give the costs $\tau_G$ of the time-based allocation. The sample count of every unit is verified, uploads
are read back, and every wait on a device is guarded by a watchdog; a device that fails is removed and its units are
requeued on the others. Since a unit's result depends only on its counters, the result does not depend on how the work is
split among the devices. The kernel's limits are compile-time constants, so the larger graphs need larger builds of the
same sources: the default one handles up to 16 lines and 768 atoms per class, enough through eight loops; a variant with
up to 18 lines and 8192 atoms covers nine loops (17 lines for $\Gamma^{(2)}$, 18 for $\Gamma^{(4)}$), and one with up to
20 lines, subgraphs of up to 11 loops and 16\,384 atoms covers ten. The sector table of a class has $2^\ell$ entries,
about a million at ten loops. Each build has its own signature, and the ledger accepts runs only from known
ones (Sec.~\ref{sec:methods-ledger}).

% ---------------------------------------------------------------------------------------------------------------
\subsection{Resummation}
\label{sec:methods-resummation}

The critical exponents follow from the RG functions at the non-trivial zero $\ubar^*$ of $\beta$: $\eta=\eta(\ubar^*)$, $\nu=\nu(\ubar^*)$ and $\omega=\beta'(\ubar^*)$. The series are asymptotic, with
coefficients that grow like $k!$ at large order $k$~\cite{Lipatov1977,BrezinLeGuillouZinnJustin1977}, and $\ubar^*$
is of order one, so they must be resummed. The eighth-order series of Ref.~\cite{Phi4tools2024} were not resummed
there (``We postpone such study to a future work''), and to our knowledge they have not been resummed since; the
standard fixed-dimension estimates rest on the shorter series analysed in
Refs.~\cite{LeGuillouZinnJustin1977,LeGuillouZinnJustin1980,GuidaZinnJustin1998}. We resum the eighth-order series,
from Ref.~\cite{Phi4tools2024} and from our own coefficients, and our ninth- and tenth-order series with the same
method and the same settings, so that the effect of the new terms, and that of the more precise eighth-order
coefficient of $\eta$, can be separated from that of the method.

\subsubsection{Borel--Leroy transform and conformal mapping}
\label{sec:resummation-method}

We follow Refs.~\cite{LeGuillouZinnJustin1977,LeGuillouZinnJustin1980,GuidaZinnJustin1998} (see also
Ref.~\cite{ZinnJustinBook}). For a series $F(\ubar)=\sum_k F_k\,\ubar^k$ the Borel--Leroy transform
\begin{gather}
F(\ubar)=\int_0^\infty dt\;t^{b}\,e^{-t}\,\mathcal B_F^{(b)}(\ubar t),\nonumber\\
\mathcal B_F^{(b)}(x)=\sum_k\frac{F_k}{\Gamma(k+b+1)}\,x^k,
\label{eq:borel}
\end{gather}
has a convergent expansion, since at large order~\cite{Lipatov1977,BrezinLeGuillouZinnJustin1977,BrezinParisi1978,ZinnJustinBook}
\begin{equation}
F_k\simeq c_F\,k!\,(-a)^k\,k^{b_F}\bigl[1+O(1/k)\bigr].
\label{eq:largeorder}
\end{equation}
The constant $a=48\pi/[(N+8)A]$, fixed by the action $A=\rsInstanton$ of the instanton
(Sec.~\ref{sec:methods-validation}), is the same for all RG functions ($a=\rsAone$ at $N=1$), while $b_F$ and $c_F$
depend on $F$: $b_F=3+N/2$ for $\beta$, for $\eta_{\phi^2}$ and hence for $\nu$ and $1/\nu$, and $b_F=2+N/2$ for
$\eta$~\cite{ZinnJustinBook,Ribeca2003thesis}. Assuming that $\mathcal B_F^{(b)}$ is analytic in the plane cut along
$(-\infty,-1/a]$, the mapping
\begin{equation}
w(x)=\frac{\sqrt{1+a x}-1}{\sqrt{1+a x}+1}
\label{eq:conformal}
\end{equation}
sends the cut plane onto the unit disk, and the series in $\ubar$ is replaced by
\begin{equation}
F(\ubar)\approx\sum_{k=0}^{K}d_k(\alpha,b)\int_0^\infty dt\;t^{b}\,e^{-t}\,
\frac{[w(\ubar t)]^k}{[1-w(\ubar t)]^{\alpha}},
\label{eq:resummed}
\end{equation}
where the $d_k(\alpha,b)$ are fixed by matching the expansion of the right-hand side to $F_0,\dots,F_K$. The factor
$[1-w]^{-\alpha}$ makes the approximant of the Borel transform grow like $x^{\alpha/2}$ at large $x$ (it is the
parameter $\lambda=\alpha/2$ of Ref.~\cite{KompanietsPanzer2017}). Dividing $F$ by $\ubar^s$ is equivalent to replacing
$(b,\alpha)$ with $(b-s,\alpha+2s)$, so the scan over $\alpha$ below contains the shifted series of
Ref.~\cite{GuidaZinnJustin1998}, which also used a homographic transformation of the coupling and a factor
$1+r\ubar$ that we do not use.

\subsubsection{Resummed functions, parameters and errors}
\label{sec:resummation-criteria}

The fixed point $\ubar^*$ is the zero of the resummed $\beta$, and $\omega$ is the slope of the resummed $\beta$ there;
$\eta$ is resummed and evaluated at $\ubar^*$. For $\nu$ we resum the series of $1/\nu=2-\eta+\eta_{\phi^2}$ as a
single series and take its inverse: the Callan--Symanzik relations~\eqref{eq:csrel} give $1/\nu$ directly, as a
product of series, whereas the series of $\nu$ follows from it by an inversion, which turns any zero of $1/\nu$ in the
complex $\ubar$ plane into a pole. As Ref.~\cite{GuidaZinnJustin1998} summed both $S$ and $1/S$ for each exponent,
we also resum the series of $\nu$ itself as a second determination; the two differ by at most $\rsNuDirMax$, and the
difference enters the error~\check{S2}.

We scan $b=0,\frac12,\dots,\rsBmax$ and $\alpha=\rsAlphaMin,\dots,\rsAlphaMax$ in steps of $\frac12$; the range of $b$
contains the large-order values $b_F$, and beyond $\alpha\approx\rsAlphaMax$ the successive truncations oscillate with
growing amplitude (those of $\eta$ from $\alpha\approx12$; see Sec.~\ref{sec:resummation-results}). For each
$(b,\alpha)$ we compute the results $F^{[K-2]}$, $F^{[K-1]}$ and $F^{[K]}$ of the last three truncations of
Eq.~\eqref{eq:resummed} and, over the neighbourhood $|b'-b|\le1$, $|\alpha'-\alpha|\le1$,
the largest change $\delta$ between successive truncations and half the range $\sigma$ of $F^{[K]}$. As in
Ref.~\cite{GuidaZinnJustin1998}, the central parameters $(b^*,\alpha^*)$ are those for which the result depends least
on the truncation order and on the parameters, here the minimum of $\delta+\sigma$; taking the largest change over a
neighbourhood rather than at the point itself discards the isolated points where successive truncations cross by
accident, which are frequent at large $\alpha$. The region of comparable stability is the set $R$ of grid points
where $\delta+\sigma$ is at most twice its minimum. The resummation error at fixed $\ubar^*$ is the largest of
$\delta$, $\sigma$ and $\max_R|F^{[K]}-F^{[K]}(b^*,\alpha^*)|$ (for $\nu$ also the difference from the series of $\nu$); as in
Ref.~\cite{GuidaZinnJustin1998}, the change of the result when $\ubar^*$ moves by its own error is added to it. Every
series is treated with the same grid and the same criteria. The truncation order is $K=8$ for the eighth-order series,
$K=9$ for the ninth-order ones and, at tenth order, $K=10$ for $\eta$ and $1/\nu$ and $K=11$ for $\beta$.

\subsubsection{Propagating the Monte Carlo uncertainties}
\label{sec:resummation-errors}

The coefficients carry statistical errors, which we propagate through the resummation and keep separate from its
systematic error. Every per-order sum is stored class by class (Appendix~\ref{app:data}), with independent random
numbers for every class and loop order (Sec.~\ref{sec:methods-rng}), and is linear in the class integrals: if these
are drawn from Gaussians with their quoted errors, each sum at fixed $N$ is a Gaussian variable with the propagated
error of Sec.~\ref{sec:results-new}, independent of the other sums. We therefore draw the sums at each $N$ (for the
orders taken from Ref.~\cite{Phi4tools2024}, and for its own coefficients in the eighth-order row of
Table~\ref{tab:exponents}, with their quoted errors), rebuild the coefficients through Eq.~\eqref{eq:csrel}, and resum
every draw with the central parameters; the standard deviation over \rsDraws\ draws, with a fixed seed, is the
statistical error. The sums at different $N$ are correlated through the shared class integrals, but no entry of
Table~\ref{tab:exponents} combines different $N$. Drawing the class integrals themselves, for the ledger inputs of the
eighth-order series at $N=1$, gives the same spreads within \rsClassDev, as expected from the sampling error of
\rsDraws\ draws, and a linearized propagation agrees with the draws within \rsLinDev~\check{S2}.

The statistical errors of our results are at most $\rsStatMaxUstar$ for $\ubar^*$, $\rsStatMaxNu$ for
$\nu$, $\rsStatMaxEta$ for $\eta$ and $\rsStatMaxOmega$ for $\omega$, at most $\rsStatRatio$ times the resummation
errors (the largest ratios, for $\ubar^*$ at $N=2$ and~3, come mostly from the coefficient of $\ubar^{11}$ in
$\beta$). With the quoted errors of the eighth-order coefficients of Ref.~\cite{Phi4tools2024}, larger than ours, the
ratio is at most $\rsStatRatioA$~\check{S2}. The resummed exponents are therefore limited by the resummation, not
by the precision of the coefficients: more computing time spent on the present orders would not change them. Further
progress needs more orders, which have so far reduced the errors only gradually (Table~\ref{tab:exponents}), and a
better control of the non-analyticity at $\ubar^*$ (Sec.~\ref{sec:resummation-results}).

% ---------------------------------------------------------------------------------------------------------------
\subsection{Validation strategy}
\label{sec:methods-validation}

Individual subtracted integrals have no external reference; only the order sums do. The checks are therefore organized as
follows; their results are given in Sec.~\ref{sec:results}.

\paragraph{Integrands.}
The expansion~\eqref{eq:integrand} agrees pointwise with an independent implementation in Python (to
$1.4\times10^{-14}$ of the sum of the absolute values of the terms, on 39 graphs at two to five loops) and with an
independent multiprecision implementation with its own forests, Kirchhoff polynomials from the matrix-tree theorem and no
pruning (all graphs at three to five loops and 20 graphs at six loops); pruned and unpruned expansions agree to the
round-off of extended precision~\check{I1}. The vanishing graphs, the twins
(Sec.~\ref{sec:methods-graphs})~\checks{I2,I3} and the shifted faces
(Sec.~\ref{sec:methods-tropical}) are checked as described above.

\paragraph{Exact values and known orders.}
The exact values of $B_1$, $B_2$, $H_2$, $C_1$ and $C_2$, and that of the subtracted integral of one three-loop graph (a
sunset inserted in a ring), are reproduced at high statistics. The coefficients known from
Ref.~\cite{Phi4tools2024}, $B_L$ to $L=7$ and $H_L$, $C_L$ to $L=8$, follow from its series through
Eqs.~\eqref{eq:csrel} and~\eqref{eq:series}, with errors propagated from those quoted there~\check{S1}; a run is compared with them
through the pull
\begin{equation}
z=\frac{X-X_{\mathrm{ref}}}{\sqrt{\delta X^2+\delta X_{\mathrm{ref}}^2}} ,
\label{eq:pull}
\end{equation}
where $X\pm\delta X$ is our value of an order sum and $X_{\mathrm{ref}}\pm\delta X_{\mathrm{ref}}$ the reference, at
$N=1$ and, by recombining the same integrals with their $O(N)$ polynomials, at $N=0$, 2 and 3. The unsubtracted intermediate-scheme integrals, computed by the
same code, are compared graph by graph with the values of Ref.~\cite{Phi4tools2024} and with their order sums.

\paragraph{Error bars.}
Replica ensembles with independent seeds test the error estimates: the mean and spread of their pulls, and the scatter of
every class among replicas against its quoted error ($\chi^2$ per degree of freedom). The same per-class $\chi^2$ over
repeated runs tests classes with large kurtosis or large maximal weights individually. Since every class's integral is
stored, two runs can be compared class by class, independently of any reference.

\paragraph{Hardware.}
Golden runs and device self-tests (Sec.~\ref{sec:methods-kernel}) establish that the GPUs compute the same numbers as the
CPU, up to the round-off of the device library, before any physics run is accepted~\check{N2}.

\paragraph{New orders.}
For $B_8$--$B_{10}$, $H_9$, $H_{10}$, $C_9$ and $C_{10}$, which have no reference, and for $H_8$ and $C_8$, whose
precision we improve, we rely on the following checks.
\begin{enumerate}
\item \emph{A second route} (route B). The coefficients are also obtained from the intermediate-scheme
integrals of $\Gamma^{(4)}(0)$, $\partial\Gamma^{(2)}/\partial p^2(0)$ and $\Gamma^{(2)}(0)$, in which only the sunset
subgraphs are subtracted, through the $Z$ factors and the reversion of the coupling series~\cite{Ribeca2003thesis}, as in
Ref.~\cite{Phi4tools2024}. In this scheme $\Gamma^{(2,1)}(0;0)$ follows from $\Gamma^{(2)}(0)$ through the exact
identity~\cite{Ribeca2003thesis}
\begin{equation}
\bigl(2-u_I\,\partial_{u_I}\bigr)\hat\Gamma^{(2)}_I(0)
=\Bigl(2+\tfrac{N+2}{3}\,u_I-\tfrac{2(N+2)}{9}\,u_I^2\Bigr)\Gamma^{(2,1)}_I(0;0),
\label{eq:intermediate}
\end{equation}
with $u_I=g_I/8\pi$, where $g_I=\lambda_0/m_I$ is the bare coupling in units of the mass $m_I$ of the intermediate
scheme, $\hat\Gamma^{(2)}_I=\Gamma^{(2)}_I/m_I^2$ and integrals in Nickel's normalization; the same identity
checks intermediate-scheme data order by order. Route B shares with route A only the graph lists, the weights and the sampler, and,
because of the cancellations of Sec.~\ref{sec:methods-cs}, it reaches a lower precision at equal cost.
The formulas are derived to all orders and reproduce every published coefficient of Ref.~\cite{Phi4tools2024}
through eighth order from its per-graph values, at $N=0$--$4$~\check{S4}. We evaluated the intermediate-scheme
integrals with eight to ten vertices on GPUs, in \routeBhours\ GPU hours, and took those with fewer vertices from
Ref.~\cite{Phi4tools2024}; the ninth- and tenth-order coefficients of the two routes agree within \routeBzmax\
standard deviations (Sec.~\ref{sec:results-new}, Table~\ref{tab:routes})~\check{S4}. Route B resolves them only to
\routeBrelmin--\routeBrelmax\ relative, and it has no check of the coefficient of $\ubar^{11}$ in $\beta$, which would
need the four-point graphs with eleven vertices.
\item \emph{Sum rules}: the eight- to ten-loop lists and weights pass the zero-dimensional sum rules
(Sec.~\ref{sec:methods-graphs})~\check{G2}.
\item \emph{Large $N$}: the leading power of $N$ of $B_L(N)$, $H_L(N)$ and $C_L(N)$ at each order is compared with the
$1/N$ expansion in the same scheme~\cite{ZinnJustinBook,Ribeca2003thesis}. At $N=\infty$ with $\ubar$ fixed only the
bubble chains survive, and $\beta=-\ubar+\ubar^2$, $\eta=0$, $\eta_{\phi^2}=-\ubar$, $\nu=1/(2-\ubar)$: the terms in
$N^L$ of $B_L$, $H_L$ and $C_L$ must vanish for $L\ge2$. No class of the ledger has an $O(N)$ weight of degree $L$, so
they vanish identically, and in the normalization of Eq.~\eqref{eq:series-form} (Sec.~\ref{sec:results-new})
$p_{\beta,k}$ has degree $k-2$, $p_{\eta,k}$ and
$p_{\eta_{\phi^2},k}$ degree $k-1$, and the coefficient of $N^k$ in $p_{\nu,k}$ is $2^{-k-1}$. The terms in $N^{L-1}$
follow from the $1/N$ corrections to $Z$, $\Gamma^{(4)}(0)$ and $\Gamma^{(2,1)}(0;0)$, one-dimensional integrals over
the momentum of the bubble chain, which reproduce the leading coefficients of the polynomials of
Ref.~\cite{Phi4tools2024} through eighth order to \largeNphi~\check{S5}.
\item \emph{Large order}: every coefficient is compared with its leading large-order estimate
$c\,(-a)^k\,\Gamma(k+b+1)$, equivalent at large $k$ to Eq.~\eqref{eq:largeorder}~\cite{Lipatov1977,BrezinLeGuillouZinnJustin1977,BrezinParisi1978,ZinnJustinBook},
with the instanton constants of $d=3$~\cite{BrezinLeGuillouZinnJustin1977,ZinnJustinBook}, $A=113.3835$ and the
amplitudes $c$, so that in Nickel's coupling $\ubar$ the constant is $a=48\pi/[(N+8)A]=0.147774232\times9/(N+8)$
[Ref.~\cite{ZinnJustinBook}, Eq.~(28.73)], and $b=3+N/2$ for $\beta$ and $\eta_{\phi^2}$, $2+N/2$ for $\eta$, as
implemented for the seven-loop computation~\cite{Ribeca2003thesis}.
The leading estimate alone catches an error of a few per cent in a new coefficient~\check{S3}.
\end{enumerate}

